\clearpage\thispagestyle{empty}
\pdfbookmark[0]{Quellen zum Brief vom 14. Januar 1811}{lindenau1811jan-sources}
\begin{center}\textbf{Quellen zum Brief vom 14. Januar 1811}\end{center}
\noindent Diese KI-Übertragung ins heutige Deutsch folgt dem Brief von Gauß an Bernhard August von Lindenau in \href{https://gdz.sub.uni-goettingen.de/id/PPN236007467}{Gauß’ \textit{Werken}, Bd.~VI} (Göttingen, 1874), S.~325. Der Erstdruck erschien als Beitrag VIII der Zeitschrift \textit{Monatliche Correspondenz}, Bd.~XXIII, Januar 1811, S.~\href{https://api.digitale-sammlungen.de/iiif/image/v2/bsb10538612_00105/full/1600,/0/default.jpg}{97} und \href{https://api.digitale-sammlungen.de/iiif/image/v2/bsb10538612_00106/full/1600,/0/default.jpg}{98}. Das digitalisierte Exemplar stammt aus der Bayerischen Staatsbibliothek.
\par\smallskip\noindent Gauß verweist auf die Elemente im \href{https://archive.org/download/bub_gb_YCI4AAAAMAAJ/page/n595.jpg}{Dezemberheft von 1810, S.~592}. Die tägliche mittlere tropische Bewegung steht in der Werkausgabe des Briefs als $768.980$. Der Erstdruck verwendet den Doppelstrich und das Dezimalkomma. Er enthält außerdem die Ephemeride vom 3. Februar bis zum 12. April 1811. In der Werkausgabe stehen dafür die Daten in eckigen Klammern. KI-Übertragung und KI-Satz.
\clearpage
